%% =====================================================================
%%  COS 802 Research Report
%%  Business-Aware Customer Segmentation: A Comparative Evaluation of
%%  Standard, RFM-Based and Constrained K-Means Clustering
%%
%%  Template : IEEEtran (conference mode, two-column)
%%  Structure: SIGKDD Explorations research-report layout
%%             (Abstract, Keywords, Introduction, Related Work, Data,
%%              Methodology, Experiments & Results, Discussion,
%%              Conclusion & Future Work, Reproducibility, References)
%%  Results  : every number in this report was produced by an independent
%%             re-run of the project notebook on the original UCI file
%%             (Online Retail.xlsx, 541,909 rows), extended with k-selection,
%%             cross-space validation, bootstrap confidence intervals,
%%             outlier-robust effect sizes and baseline constraint rates.
%%  Compile  : pdfLaTeX. Figures are pre-rendered PDFs in figures/;
%%             their pgfplots sources are in figures/src/ (compile each
%%             with pdflatex and copy the PDF up one level to change one).
%% =====================================================================
\documentclass[conference,10pt]{IEEEtran}
\IEEEoverridecommandlockouts

% ---------- Packages ----------
\usepackage[utf8]{inputenc}
\usepackage[T1]{fontenc}
\usepackage{cite}
\usepackage{amsmath,amssymb,amsfonts}
\usepackage{graphicx}
\usepackage{textcomp}
\usepackage{xcolor}
\usepackage{booktabs}
\usepackage{multirow}
\usepackage{array}
\usepackage{algorithm}
\usepackage{algpseudocode}
\usepackage{url}
\usepackage[hidelinks]{hyperref}

% ---------- Colour-blind-safe palette ----------
\definecolor{cStd}{RGB}{0,114,178}    % Standard K-means
\definecolor{cRFM}{RGB}{230,159,0}    % RFM K-means
\definecolor{cCon}{RGB}{0,158,115}    % Constrained K-means
\definecolor{cML}{RGB}{213,94,0}      % must-link
\definecolor{cCL}{RGB}{86,180,233}    % cannot-link
% segment archetypes (PCA figure)
\definecolor{aHigh}{RGB}{213,94,0}
\definecolor{aReg}{RGB}{0,114,178}
\definecolor{aOcc}{RGB}{0,158,115}
\definecolor{aLap}{RGB}{160,160,160}

\newcommand{\etasq}{\eta^{2}}
\newcolumntype{L}[1]{>{\raggedright\arraybackslash}p{#1}}

\begin{document}

% =====================================================================
\title{Customer segmentation and behavioural analysis: A Comparison of
clustering methods for e-commerce segmentation}

\author{\IEEEauthorblockN{Nqobile Ngidi}
\IEEEauthorblockA{\textit{Student number: u15189181}\\
\textit{COS 802 Research Project, Department of Computer Science}\\
\textit{University of Pretoria}, Pretoria, South Africa\\
u15189181@tuks.co.za}}

\maketitle

% =====================================================================
\begin{abstract}
Customer segmentation is a widely used approach
in the identification of customer purchase behaviours. These
segments are used to support business decisions relating to
marketing, resource allocation and product development. However, conventional customer segmentation does not explicitly
reflect some business metrics. This study compares standard
K-means, K-means on Recency-Frequency-Monetary (RFM), and a
soft constrained K-means using the UCI Online Retail dataset.
Segments were built using historic transaction data and evaluated
against outcomes in a subsequent 90-day period that was reserved
as a temporal holdout. The model performance was measured
using Silhouette, Davies–Bouldin and Calinski–Harabasz indices,
and eta-squared measured the future business-metric separation; constraint violation rates and sensitivity analyses were also reported. The
results show that RFM K-means achieved the best internal
scores (Silhouette 0.370) and the largest future-order separation
($\eta^2 = 0.192$ versus 0.147; bootstrap 95\% CI of the difference
[0.020, 0.086]). Constrained K-means changed the assignment
of only 48 customers (1.4\%) relative to standard K-means,
raised future-order $\eta^2$ by 0.006 and lowered it in two of five
resampled constraint sets, exhibiting only marginal improvements.
The study demonstrates the importance of evaluating customer
segmentation through both statistical validity and future business
relevance. It concludes that, within the experimental setting,
RFM-based K-means provides the strongest overall performance,
while the tested country-based constraints do not establish
a meaningful advantage over standard K-means. The study
contributes a reproducible, temporally validated framework that
evaluates validity, business-oriented metrics and constraint
sensitivity.
\end{abstract}

\begin{IEEEkeywords}
customer segmentation, K-means, RFM analysis, constrained clustering,
must-link and cannot-link constraints, temporal holdout, cluster validation,
bootstrap, e-commerce analytics
\end{IEEEkeywords}

% =====================================================================
\section{Introduction}
\label{sec:intro}

\subsection{Context and motivation}
The growth of e-commerce sales has brought with it a significant expansion in e-commerce retailer data, which includes information such as customer transaction date, product type, quantity, price and unique customer identifier. These features can be used to identify customer purchase behaviours and segments, describing how often customers purchase and an estimated purchase spend. Customer segmentation has the ability to
group customers with similar behaviour so that communication, retention
effort and service levels can be tailored to each group
\cite{chen2012rfm,salminen2023algorithms}. Because transactional datasets rarely
contain segment labels, segmentation is usually framed as unsupervised
clustering \cite{wani2024comprehensive,gomes2023review}.

K-means is the most widely used model for customer segmentation; it is advantageous because it is simple to compute and
fast
\cite{macqueen1967some,gomes2023review,salminen2023algorithms}. In e-commerce it is commonly used with
the RFM framework \cite{chen2012rfm,christy2021rfm}. A further
option is constrained (semi-supervised) clustering, in which domain knowledge
is injected as pairwise must-link and cannot-link constraints
\cite{wagstaff2001constrained,basu2004probabilistic}. However, internal clustering and business relevance may not always be equivalent. A model may effectively segment customers according to recency, frequency and monetary spend but may overlook other segments that are informed by business such as geographic location.

\subsection{Problem statement}
Past studies on segmentation typically evaluate clustering by internal indices
such as the Silhouette score \cite{rousseeuw1987silhouettes}. These metrics indicate how well separated groups are in the clustering feature space, but they do not speak to separation in terms of behaviour that matters to the
business \cite{salminen2023algorithms}. Constrained
clustering is an approach that can be used to make segments ``business-aware''. However, its benefit
depends on the quality of the constraints and on how strongly they are enforced. 

Clustering methods should not only be evaluated on internal cluster separation but should be evaluated in relation to business-relevant behaviour. 

\subsection{Research questions}
\textbf{RQ1.} How does the performance of standard K-means clustering
statistically compare with RFM-based K-means and constrained K-means
clustering in customer segmentation?

\textbf{RQ2.} How does constrained clustering measurably improve
business-metric separation relative to unconstrained clustering?

\textbf{RQ3.} is analysed as the difference in $\etasq$ between constrained and
standard K-means, which share a feature space. This makes the comparison
measurable.

\subsection{Contributions}
\begin{enumerate}
  \item A controlled comparison of standard, RFM and soft constrained K-means
  at a common $k$.
  \item A temporal-holdout evaluation of five future outcomes.
  \item An explicit audit of the constrained model: violation rates compared
  with unconstrained and chance baselines and penalty-weight sensitivity.
\end{enumerate}

Section~\ref{sec:related} reviews related work, Section~\ref{sec:data}
describes the data and Section~\ref{sec:method} the methodology.
Section~\ref{sec:results} reports the experiments, Section~\ref{sec:discussion}
discusses them, and Section~\ref{sec:conclusion} concludes.

% =====================================================================
\section{Related Work}
\label{sec:related}

\subsection{RFM-based customer segmentation}
The RFM framework represents customer behaviour through Recency which measures the time since the last purchase, Frequency which measures the number of purchases within a period, lastly, Monetary value which measures spending value. The framework originates from marketing analytics and
it is the most common representation for transaction-based segmentation
\cite{gomes2023review}. Chen et~al.\
\cite{chen2012rfm} applied RFM with K-means to the 
UK online gift retail data and profiled the clusters into actionable
groups. Christy et~al.\ \cite{christy2021rfm} proposed an RFM ranking scheme
and compared K-means with fuzzy C-means and a modified K-means. The
\emph{merit} of this work is interpretability and low cost; the \emph{gap}
is that RFM omits basket size, product breadth and purchase regularity, and
that segments are usually evaluated on the same variables used to form
them. These studies on the RFM based k-means form a basis for this study.

\subsection{Comparative studies of clustering algorithms}
K-means, introduced by MacQueen \cite{macqueen1967some}, is an algorithm that partitions data observations into clusters using various iterations to assign observations to its nearest centroid. This approach is valued for its
scalability and interpretability, but it is
sensitive to scale and outliers.  John et~al.\
\cite{john2023exploration} used the UCI  retail dataset and compared K-means, a Gaussian mixture model (GMM),
DBSCAN \cite{ester1996dbscan}, agglomerative clustering and BIRCH and found
that the GMM performed best (Silhouette 0.80). At a larger scale, Salminen
et~al.\ \cite{salminen2023algorithms} reviewed 172 segmentation studies and
found 46 different algorithms and 14 evaluation metrics, with K-means being the most common algorithm. These
studies motivate controlled comparisons but rank methods
predominantly by internal clustering indices.

\subsection{Internal cluster validation}
The Silhouette coefficient \cite{rousseeuw1987silhouettes} measures and contrasts the closeness of a point to it's cluster with its distance to the nearest other
cluster; the Davies--Bouldin measures the within cluster dispersion relative to the between cluster seperation; and the Calinski--Harabasz index is a
variance-ratio criterion \cite{ikotun2025cluster}. All three are defined relative to a feature space, so they are not
directly comparable across representations and do not measure business
usefulness \cite{salminen2023algorithms}.

\subsection{Constrained and semi-supervised clustering}
Wagstaff et~al.\ \cite{wagstaff2001constrained} introduced COP-KMeans, which
enforces must-link and cannot-link constraints as hard rules. Davidson and
Ravi \cite{davidson2005clustering} discuss that hard enforcement of cannot link constraints is NP- complete. Thus,leading the to soft constraints, soft
formulations such as PCKMeans \cite{basu2004probabilistic} add a penalty for
each violated constraint to the K-means objective. A study by Gonz\'alez-Almagro et~al.\ \cite{gonzalez2025semi}, analyses
315 constrained clustering methods, performed on pairwise instance-level
constraints, and ranks the methods by popularity and by the quality of their
experimental validation. The \emph{merit} is that
limited domain knowledge can steer clusters toward useful partitions; the
\emph{gap} in customer analytics is that public retail data rarely contain
such knowledge, and studies seldom report how many constraints are satisfied
relative to an unconstrained baseline, or how sensitive results are to the
penalty weight and to the sampled pairs.

\subsection{Business-oriented evaluation of segments}
Alves Gomes and Meisen \cite{gomes2023review} reviewed 105 e-commerce
segmentation studies published between 2000 and 2022 and found that customers
are mostly represented by hand-picked features or RFM variables. Salminen
et~al.\ \cite{salminen2023algorithms} add that separation-focused metrics
dominate evaluation and that domain experts are rarely involved.
Evaluating segments on a later period, and quantifying the uncertainty of
that evaluation, is uncommon in the clustering literature.

\subsection{Summary and research gap}
Table~\ref{tab:litgap} summarises the reviewed work. The gap filled here is a
\emph{controlled, temporally validated comparison} of standard, RFM and
constrained K-means that reports (a) internal validity in own and common
spaces, (b) separation of future outcomes not used in clustering, with
uncertainty, and (c) constraint satisfaction. 

\begin{table}[t]
\centering
\caption{Merits and gaps of related work and how this study responds.}
\label{tab:litgap}
\footnotesize
\setlength{\tabcolsep}{3pt}
\begin{tabular}{L{1.55cm}L{2.05cm}L{2.05cm}L{1.85cm}}
\toprule
\textbf{Work} & \textbf{Merit} & \textbf{Gap} & \textbf{This study}\\
\midrule
Chen et~al.\ \cite{chen2012rfm}, Christy et~al.\ \cite{christy2021rfm}
& Interpretable RFM segments on retail data
& Evaluated on same RFM variables; RFM only
& Adds extended features and a future holdout\\
\addlinespace
John et~al.\ \cite{john2023exploration}
& Multi-algorithm comparison on UK retail
& Internal indices only; no business outcomes
& Adds $\etasq$ with CIs on five future outcomes\\
\addlinespace
Wagstaff et~al.\ \cite{wagstaff2001constrained}, Basu et~al.\ \cite{basu2004probabilistic}
& Constraints inject domain knowledge
& Assume reliable constraints; little use in retail
& Audits proxy constraints against baselines\\
\addlinespace
Alves Gomes \& Meisen \cite{gomes2023review}, Salminen et~al.\ \cite{salminen2023algorithms}
& Systematic reviews of segmentation practice
& K-means/RFM dominate; business evaluation rare
& Adds holdout business outcomes\\
\addlinespace
Gonz\'alez-Almagro et~al.\ \cite{gonzalez2025semi}
& Survey of 315 constrained methods
& Few retail applications
& Applies soft constraints to retail data\\
\bottomrule
\end{tabular}
\end{table}

% =====================================================================
\section{Data}
\label{sec:data}

\subsection{Dataset}
The UCI \emph{Online Retail} dataset \cite{uci2015onlineretail} contains
541,909 transaction lines from a UK-based, non-store online gift retailer
between 1~December~2010 and 9~December~2011. The transaction data includes information such as invoice number, product
code, quantity, timestamp, unit price, customer identifier and
country.

\subsection{Cleaning}
Table~\ref{tab:cleaning} shows the effect of the cleaning proces on the data size. The cleaning process removed lines without a
customer identifier (24.9\% of rows) ; canceled invoices and lines with non-positive
quantity or price were removed. The revenue feature was computed as quantity
$\times$ unit price; revenue is therefore \emph{observed gross transaction
revenue}, not accounting revenue or profit.

\begin{table}[t]
\centering
\caption{Data cleaning flow.}
\label{tab:cleaning}
\footnotesize
\setlength{\tabcolsep}{3pt}
\begin{tabular}{L{4.6cm}rr}
\toprule
\textbf{Step} & \textbf{Rows} & \textbf{Cust.}\\
\midrule
Raw file                                   & 541,909 & 4,372\\
Drop rows with missing fields (135,080 lack ID) & 406,829 & 4,372\\
Remove 8,905 cancellation lines and 40 non-positive lines & 397,884 & 4,338\\
\midrule
Training window ($<$ 10 Sep 2011)          & 236,358 & 3,370\\
Future holdout ($\geq$ 10 Sep 2011)         & 161,526 & 1,921$^{\dagger}$\\
\bottomrule
\multicolumn{3}{l}{\scriptsize $^{\dagger}$Returning customers from the training population.}
\end{tabular}
\end{table}

\subsection{Temporal split}
The holdout begins 90 days before the last observed date (cut-off
10~September~2011). The training window (1~Dec~2010 08:26 to 9~Sep~2011
15:53) is the only data used for features, constraints and clusters; the
holdout (11~Sep~2011 10:35 to 9~Dec~2011 12:50) is used only for model evaluation.
There were 3,370 unique customers in the training observations and 1,449 (43.0\%) were do not appear in the future outcome. Customers first seen in the
holdout are excluded because they have no training features.

\subsection{Data characteristics}
90.1\% of the customers are British as per the identified modal country (36 countries in total).
All behavioural variables are strongly right-skewed (Table~\ref{tab:desc}).
The top 1\% of customers (34) generate 38.7\% of future revenue. 

\begin{table}[t]
\centering
\caption{Customer-level descriptive statistics ($n=3{,}370$).}
\label{tab:desc}
\footnotesize
\setlength{\tabcolsep}{3.5pt}
\begin{tabular}{lrrrr}
\toprule
\textbf{Variable} & \textbf{Median} & \textbf{Mean} & \textbf{Max} & \textbf{Skew}\\
\midrule
Recency (days)        & 74     & 95.0    & 283      & 0.8\\
Frequency (invoices)  & 2      & 3.5     & 132      & 9.3\\
Monetary (\pounds)    & 566    & 1,619   & 178,303  & 16.1\\
Avg.\ order value (\pounds) & 284 & 398  & 77,184   & 47.4\\
Total units           & 314    & 952     & 127,744  & 18.1\\
Product diversity     & 29     & 49.9    & 1,230    & 6.2\\
Mean purchase gap$^{*}$ & 52.2 & 64.7    & 271      & 1.5\\
\midrule
Future revenue (\pounds) & 174 & 846     & 168,470  & 22.5\\
Future orders         & 1      & 1.50    & 89       & 12.9\\
Future AOV (\pounds)  & 141    & 294     & 168,470  & 56.5\\
\bottomrule
\multicolumn{5}{l}{\scriptsize $^{*}$Computed for the 1,956 customers with $\geq 2$ purchase dates.}
\end{tabular}
\end{table}

% =====================================================================
\section{Methodology}
\label{sec:method}

\subsection{Research design}
The studyis experimental workflow can be seen on  (Fig.~\ref{fig:pipeline}). Standard and constrained K-means
share the extended feature space. RFM K-means uses a narrower representation, so its comparison
reflects representation as well as algorithm.

\begin{figure}[t]
\centering
\includegraphics{figures/fig_pipeline.pdf}
\caption{Experimental workflow. Future outcomes are never used to build
features, constraints or clusters.}
\label{fig:pipeline}
\end{figure}

\subsection{Feature engineering}
Table~\ref{tab:features} defines the seven training-window features. 

\begin{table}[t]
\centering
\caption{Customer-level features computed from the training window.}
\label{tab:features}
\footnotesize
\setlength{\tabcolsep}{3pt}
\begin{tabular}{L{2.2cm}L{4.4cm}c}
\toprule
\textbf{Feature} & \textbf{Definition} & \textbf{Set}\\
\midrule
Recency & Days from last purchase to snapshot date & R, E\\
Frequency & Number of distinct invoices & R, E\\
Monetary & Sum of line revenue & R, E\\
Avg.\ order value & Mean revenue per invoice & E\\
Total units & Sum of quantities purchased & E\\
Product diversity & Number of distinct stock codes & E\\
Mean purchase gap & Mean days between distinct purchase dates (median-imputed) & E\\
\bottomrule
\multicolumn{3}{l}{\scriptsize R = RFM feature set; E = extended feature set.}
\end{tabular}
\end{table}



\subsection{Clustering methods}
\subsubsection{Standard and RFM K-means}

Both baselines use scikit-learn \cite{pedregosa2011scikit} with k-means++
seeding \cite{arthur2007kmeanspp}, $k=4$, 20 initialisations and seed 42;
standard K-means is fitted on $\mathbf{X}_{\text{EXT}}$ (extended feature space) and RFM K-means on
$\mathbf{X}_{\text{RFM}}$.

\subsubsection{Soft constrained K-means}
 the constrained model starts from a K-means
model and sequentially reassigns points (in random order) to the cluster
minimising squared distance plus the penalty for that point's violated
constraints, updating centroids at each step and repairing empty clusters.
The best of four initialisations (maximum 35 iterations)
is kept. Constraints are soft, so violation rates are reported as outcomes.


\subsection{Constraint generation}
No segment labels or expert pairwise judgements exist for this dataset, so
each customer's \emph{modal country} is used as a weak
proxy. Random pairs are drawn using same-country pairs as must-link and
different-country pairs cannot-link until 150 of each are collected (seed
42). 

\subsection{Evaluation measures}
\subsubsection{Internal validity (RQ1)}
For point $i$ with mean intra-cluster distance $a(i)$ and mean distance to
the nearest other cluster $b(i)$, the Silhouette (S) score is measured
\cite{rousseeuw1987silhouettes}. 
The Davies-Bouldin index (DB) measures the cluster scatter
 and the Calinski-Harabasz (CH) index \cite{ikotun2025cluster} measures
dispersion matrices. A higher Silhouette score and Calinski-Harabasz index are preferred whereas a lower Davies-Bouldin are preferred. 

\subsubsection{Future business separation (RQ2)}
The correlation ratio $\etasq$ is calculated
\cite{lakens2013calculating}. It measures the the share of variance associated with cluster assignment. It is
computed for future revenue, orders, Average Order Value (AOV), units and product diversity. Uncertainty is quantified with a paired
nonparametric bootstrap \cite{rousselet2021percentile}: customers are
resampled 1,000 times with labels held fixed and the 2.5th and 97.5th
percentiles of each $\etasq$ difference form a 95\% interval. All of these
measures are descriptive associations.

\subsection{Sensitivity and robustness}
The constrained model is refitted with $w\in\{0.1, 0.5, 1, 2,
5\}$ (three initialisations, 30 iterations), and constraint sets are
re-sampled with seeds 11, 22, 33, 44 and 55 at $w=1$.


% =====================================================================
\section{Experiments and Results}
\label{sec:results}

\subsection{Experimental setup}
The code was run on the original UCI file with
Python 3.13, NumPy 2.5, pandas 3.0, scikit-learn 1.9 and SciPy 1.18. All randomness uses
fixed seeds (42 for models, 2026 for the bootstrap).

\subsection{RQ1: Internal statistical validity}
\label{sec:rq1}

\begin{table}[t]
\centering
\caption{Internal validation at $k=4$ in each method's own feature space
(best in bold).}
\label{tab:internal}
\footnotesize
\setlength{\tabcolsep}{3.2pt}
\begin{tabular}{lccccc}
\toprule
\textbf{Method} & \textbf{$S$\,$\uparrow$} & \textbf{DB\,$\downarrow$} &
\textbf{CH\,$\uparrow$} & \textbf{Min} & \textbf{Max}\\
\midrule
Standard K-means    & 0.3303 & 1.0580 & 1,952.5 & 364 & 1,415\\
RFM K-means         & \textbf{0.3699} & \textbf{0.9705} & \textbf{2,633.6} & 473 & 1,408\\
Constrained K-means & 0.3307 & 1.0590 & 1,949.9 & 349 & 1,428\\
\bottomrule
\end{tabular}
\end{table}


\textbf{Index ranking.} In own-space scoring (Table~\ref{tab:internal}) RFM K-means is best on all three indices: Silhouette
+0.040, DB $-0.088$ and CH +35\% relative to standard K-means. 

\textbf{Effect of constraints.} after optimally matching cluster labels only
48 customers (1.4\%) are assigned differently and the within-cluster sum of squares increases only from
11,101.9 to 11,111.0 (+0.08\%). The proxy constraints therefore have no
material effect on internal validity. 


\subsection{RQ2: Future business-metric separation}

\begin{table}[t]
\centering
\caption{Raw $\etasq$ for 90-day future outcomes (best in bold). All
Kruskal--Wallis tests give $p<10^{-100}$.}
\label{tab:eta}
\footnotesize
\setlength{\tabcolsep}{4pt}
\begin{tabular}{lccc}
\toprule
\textbf{Future outcome} & \textbf{Standard} & \textbf{RFM} & \textbf{Constrained}\\
\midrule
Revenue           & 0.0375 & \textbf{0.0405} & 0.0393\\
Order count       & 0.1471 & \textbf{0.1923} & 0.1529\\
Avg.\ order value & 0.00064 & \textbf{0.00137} & 0.00063\\
Units purchased   & 0.0470 & \textbf{0.0515} & 0.0492\\
Product diversity & 0.0829 & \textbf{0.1061} & 0.0846\\
\bottomrule
\end{tabular}
\end{table}



\begin{table}[t]
\centering
\caption{Outlier-robust effect sizes: $\etasq$ on $\log(1+y)$ and
rank-based $\varepsilon^2$ (Std / RFM / Con).}
\label{tab:robusteta}
\footnotesize
\setlength{\tabcolsep}{3.5pt}
\begin{tabular}{lcc}
\toprule
\textbf{Outcome} & $\boldsymbol{\etasq_{\log}}$ & $\boldsymbol{\varepsilon^2}$\\
\midrule
Revenue   & 0.155 / \textbf{0.227} / 0.158 & 0.190 / \textbf{0.277} / 0.193\\
Orders    & 0.199 / \textbf{0.304} / 0.205 & 0.172 / \textbf{0.263} / 0.176\\
AOV       & 0.129 / \textbf{0.180} / 0.132 & 0.146 / \textbf{0.199} / 0.149\\
Units     & 0.157 / \textbf{0.229} / 0.161 & 0.182 / \textbf{0.263} / 0.186\\
Diversity & 0.145 / \textbf{0.199} / 0.148 & 0.148 / \textbf{0.201} / 0.151\\
\bottomrule
\end{tabular}
\end{table}

\begin{figure}[t]
\centering
\includegraphics{figures/fig_eta.pdf}
\caption{Share of variance in each 90-day outcome explained by segment
membership, on the raw scale (top) and the log scale (bottom).}
\label{fig:eta}
\end{figure}

\textbf{Raw separation.} RFM K-means has the largest $\etasq$ on every
outcome (Table~\ref{tab:eta}, Fig.~\ref{fig:eta}). Future order count is the
most strongly separated outcome (0.192 versus 0.147 for standard K-means).
Constrained K-means exceeds standard K-means on four outcomes by at most
0.006 and is marginally \emph{below} it on AOV.


\textbf{Robust effect sizes.} Raw $\etasq$ understates segment separation
because a handful of extreme customers dominate the total variance. On a log
scale,  segment assignment accounts for 13-30\% of the variance in
every outcome (Table~\ref{tab:robusteta}).  The ranking of methods is unchanged: RFM is best
on every outcome under both robust measures, and the constrained model adds
at most 0.006 over standard K-means. 



Fig.~\ref{fig:pca} projects the extended features onto two principal
components  \cite{jolliffe2016pca}. Standard and
constrained K-means separates the components distinctively, whereas RFM's partition are not as clear. This is as a result of it having used less features to define the customers.
Section~\ref{sec:discussion}.

\begin{figure*}[t]
\centering
\includegraphics{figures/fig_pca.pdf}
\caption{PCA projection of the extended feature space (random sample of 900
customers), coloured by segment archetype for each method.}
\label{fig:pca}
\end{figure*}

\subsection{Constraint satisfaction and penalty weight}

\begin{table}[t]
\centering
\caption{Violation rates on the 150 must-link (ML) and 150 cannot-link (CL)
country pairs, compared with unconstrained partitions and chance.}
\label{tab:baseline}
\footnotesize
\setlength{\tabcolsep}{3pt}
\begin{tabular}{L{4.2cm}cc}
\toprule
\textbf{Partition} & \textbf{ML violated} & \textbf{CL violated}\\
\midrule
Chance (random pair, standard sizes) & 69.5\% & 30.5\%\\
Standard K-means (no constraints)    & 75.3\% & 30.7\%\\
RFM K-means (no constraints)         & 70.0\% & 27.3\%\\
Constrained K-means ($w=1$)          & 72.0\% & 22.0\%\\
\bottomrule
\end{tabular}
\end{table}

Table~\ref{tab:baseline} shows that the unconstrained standard partition
violates the country constraints at almost exactly the chance rate:
behavioural clusters are unrelated to country. Adding the constraints at
$w=1$ satisfies only five more must-link pairs and 13 more cannot-link pairs
than standard K-means, and the constrained model still violates more
must-link pairs than unconstrained RFM K-means.

\begin{table}[t]
\centering
\caption{Sensitivity of constrained K-means to the penalty weight $w$.}
\label{tab:weights}
\footnotesize
\setlength{\tabcolsep}{2.6pt}
\begin{tabular}{cccccccc}
\toprule
$\boldsymbol{w}$ & \textbf{$S$} & \textbf{DB} & \textbf{ML} & \textbf{CL} &
\textbf{Rev.\ $\etasq$} & \textbf{Ord.\ $\etasq$} & \textbf{ARI}\\
\midrule
0.1 & 0.3314 & 1.058 & 74.0\% & 29.3\% & 0.0392 & 0.1527 & 0.975\\
0.5 & 0.3317 & 1.058 & 73.3\% & 26.7\% & 0.0393 & 0.1530 & 0.970\\
1.0 & 0.3307 & 1.059 & 72.0\% & 22.0\% & 0.0393 & 0.1529 & 0.963\\
2.0 & 0.3282 & 1.063 & 67.3\% & 15.3\% & 0.0393 & 0.1528 & 0.949\\
5.0 & 0.3154 & 1.084 & 48.0\% &  6.0\% & 0.0390 & 0.1510 & 0.918\\
\bottomrule
\multicolumn{8}{l}{\scriptsize ARI is agreement with the standard K-means partition.}
\end{tabular}
\end{table}


Raising $w$ from 0.1 to 5 lowers cannot-link violations from 29.3\% to 6.0\%
and must-link violations from 74.0\% to 48.0\% (Table~\ref{tab:weights}), but lowers the Silhouette from 0.331 to 0.315 and
moves the partition further from the behavioural solution (ARI 0.918).
Future revenue and order $\etasq$ do not improve; both are lowest at $w=5$.
Must-link pairs remain hard to satisfy because UK--UK pairs span the whole
behavioural range of the population, so keeping them together conflicts with
the distance term.

\subsection{Robustness to sampled constraint sets}

\begin{table}[t]
\centering
\caption{Five independently sampled constraint sets ($w=1$). Values below
the standard K-means baseline are marked $\downarrow$.}
\label{tab:robust}
\footnotesize
\setlength{\tabcolsep}{3pt}
\begin{tabular}{lcccccc}
\toprule
\textbf{Seed} & \textbf{$S$} & \textbf{Rev.\ $\etasq$} & \textbf{Ord.\ $\etasq$} & \textbf{ML} & \textbf{CL} & \textbf{ARI}\\
\midrule
11 & 0.3284 & 0.0370$\downarrow$ & 0.1450$\downarrow$ & 56.7\% & 23.3\% & 0.975\\
22 & 0.3309 & 0.0395 & 0.1531 & 64.7\% & 18.7\% & 0.964\\
33 & 0.3281 & 0.0373$\downarrow$ & 0.1467$\downarrow$ & 70.7\% & 22.7\% & 0.972\\
44 & 0.3312 & 0.0393 & 0.1533 & 68.7\% & 24.7\% & 0.964\\
55 & 0.3303 & 0.0392 & 0.1527 & 62.7\% & 18.0\% & 0.961\\
\midrule
Mean & 0.3298 & 0.0385 & 0.1502 & 64.7\% & 21.5\% & 0.967\\
SD   & 0.0014 & 0.0012 & 0.0040 & 5.5\% & 3.0\% & 0.006\\
\midrule
Standard & 0.3303 & 0.0375 & 0.1471 & -- & -- & 1\\
\bottomrule
\end{tabular}
\end{table}

Across five re-sampled constraint sets (Table~\ref{tab:robust}) the
Silhouette and $\etasq$ values are stable, but the direction of the
constrained-versus-standard difference is not: in two of five sets the
constrained model separates future revenue and orders \emph{less} well than
standard K-means. The mean gain in future-order $\etasq$ (+0.003) is smaller
than its across-set standard deviation (0.004). The improvement reported for
the primary constraint set in Table~\ref{tab:boot} is therefore specific to
that sample of pairs rather than a reliable effect of the constraints.

% =====================================================================
\section{Discussion}
\label{sec:discussion}

\subsection{Answer to RQ1}
RFM K-means produced the most compact and best separated clusters on all
three own-space indices. The internal advantage of RFM is
therefore a property of its three-dimensional representation. The direct test of constraints, standard
versus constrained K-means in the same space, shows no meaningful difference:
the two partitions agree on 98.6\% of customers. The evidence does not
support a statistical advantage for the tested constrained model.

\subsection{Answer to RQ2}
Constrained K-means did not deliver a reliable improvement in business
separation. Its gains over standard K-means were at most 0.006 in raw
$\etasq$ and 0.006 in $\etasq_{\log}$

\subsection{Why the constraints did not yield significant improvement}
Table~\ref{tab:baseline} provides insight on how unconstrained
behavioural clusters already violate the country constraints at the chance
rate, so country carries no information about purchasing behaviour in this
UK-dominated population. The result shows that using the customers' country as an informed business metric for the constraints is not effective in this context.


\subsection{Limitations}
\begin{itemize}
  \item \emph{External validity:} The data only comprised of one retailer, one year and a UK-dominated
  customer base that resulted in limited generalisation.
  \item \emph{Construct validity:} country constraints are proxies and are not informed by subject matter experts as relevent
  segment knowledge for this use case; revenue excludes returns, costs and margins; 24.9\% 
  rows of data were removed becasue they did not have a customer identifier; 42\% of customers have negative prices or quantities, which resulted in a reducuction of the feature space.
\end{itemize}

% =====================================================================
\section{Conclusion and Future Work}
\label{sec:conclusion}

\subsection{Conclusion}
This study compared standard, RFM-based and soft constrained K-means on the
UCI Online Retail dataset, combining internal validation with a temporal
holdout of future business outcomes. Three findings stand out. First, RFM
K-means scored best in interal validity and had robust gain in business separation, explaining substantially
more variance in future ordering than the extended-feature models, because
its segments are defined by recency. Third, country-based constraints changed
the assignment of 1.4\% of customers and produced business-separation gains. Segmentation should therefore be
judged on internal validity, outlier-robust separation of
independently measured outcomes, and constraint satisfaction against an
unconstrained baseline. Constraints help only when they encode significant business knowledge.

\subsection{Future work}
\begin{enumerate}
  \item Replace country proxies with behaviourally meaningful constraints.
  \item Compare the heuristic with COP-KMeans \cite{wagstaff2001constrained}
  and PCKMeans/MPCKMeans \cite{basu2004probabilistic}.
  \item evaluation over several holdout windows and a
  bootstrap that refits the clusters, so intervals include clustering
  instability.
  \item add Gaussian mixture, DBSCAN and agglomerative baselines
  \cite{john2023exploration,wani2024comprehensive}.

  \item Where data permit, include margins, returns, acquisition cost and
  campaign exposure, and validate segment-based actions with randomised
  experiments.
\end{enumerate}

% =====================================================================
\section*{Reproducibility }
The original notebook
(\texttt{Masters\_Constrained\_\allowbreak Customer\_\allowbreak Segmentation\_\allowbreak Experiment.ipynb})
and an extended script (\texttt{run\_analysis.py}) were used to inform the results in this report \texttt{Online Retail.xlsx} . The project repository is available at
\url{https://github.com/NqobileNgidi/COS802_u15189181_Customer_Segmentation}
.
and contains a README with run instructions.


% =====================================================================
\begin{thebibliography}{99}

\bibitem{chen2012rfm}
D.~Chen, S.~L. Sain, and K.~Guo, ``Data mining for the online retail
industry: A case study of RFM model-based customer segmentation using data
mining,'' \emph{J. Database Marketing \& Customer Strategy Manage.}, vol.~19,
no.~3, pp.~197--208, 2012, doi: 10.1057/dbm.2012.17. [Online]. Available:
\url{https://openresearch.lsbu.ac.uk/item/87978}

\bibitem{salminen2023algorithms}
J.~Salminen, M.~Mustak, M.~Sufyan, and B.~J. Jansen, ``How can algorithms
help in segmenting users and customers? A systematic review and research
agenda for algorithmic customer segmentation,'' \emph{J. Marketing Anal.},
vol.~11, no.~4, pp.~677--692, 2023, doi: 10.1057/s41270-023-00235-5.
[Online]. Available:
\url{https://jyx.jyu.fi/bitstream/123456789/88403/1/s41270-023-00235-5.pdf}

\bibitem{wani2024comprehensive}
A.~A. Wani, ``Comprehensive analysis of clustering algorithms: Exploring
limitations and innovative solutions,'' \emph{PeerJ Comput. Sci.}, vol.~10,
Art.~no.~e2286, 2024, doi: 10.7717/peerj-cs.2286. [Online]. Available:
\url{https://pmc.ncbi.nlm.nih.gov/articles/PMC11419652/}

\bibitem{gomes2023review}
M.~Alves Gomes and T.~Meisen, ``A review on customer segmentation methods for
personalized customer targeting in e-commerce use cases,'' \emph{Inf. Syst.
e-Bus. Manage.}, vol.~21, pp.~527--570, 2023,
doi: 10.1007/s10257-023-00640-4. [Online]. Available:
\url{https://link.springer.com/content/pdf/10.1007/s10257-023-00640-4.pdf}

\bibitem{macqueen1967some}
J.~MacQueen, ``Some methods for classification and analysis of multivariate
observations,'' in \emph{Proc. 5th Berkeley Symp. Math. Statist. Probab.},
vol.~1, 1967, pp.~281--297. [Online]. Available:
\url{https://projecteuclid.org/euclid.bsmsp/1200512992}

\bibitem{christy2021rfm}
A.~J. Christy, A.~Umamakeswari, L.~Priyatharsini, and A.~Neyaa, ``RFM
ranking---An effective approach to customer segmentation,'' \emph{J. King
Saud Univ.---Comput. Inf. Sci.}, vol.~33, no.~10, pp.~1251--1257, 2021,
doi: 10.1016/j.jksuci.2018.09.004. [Online]. Available:
\url{https://doi.org/10.1016/j.jksuci.2018.09.004}

\bibitem{wagstaff2001constrained}
K.~Wagstaff, C.~Cardie, S.~Rogers, and S.~Schr\"{o}dl, ``Constrained K-means
clustering with background knowledge,'' in \emph{Proc. 18th Int. Conf.
Machine Learning (ICML)}, 2001, pp.~577--584. [Online]. Available:
\url{https://www.wkiri.com/research/papers/wagstaff-kmeans-01.pdf}

\bibitem{basu2004probabilistic}
S.~Basu, M.~Bilenko, and R.~J. Mooney, ``A probabilistic framework for
semi-supervised clustering,'' in \emph{Proc. 10th ACM SIGKDD Int. Conf.
Knowledge Discovery and Data Mining}, 2004, pp.~59--68,
doi: 10.1145/1014052.1014062. [Online]. Available:
\url{https://www.cs.utexas.edu/~ml/papers/semi-kdd-04.pdf}

\bibitem{rousseeuw1987silhouettes}
P.~J. Rousseeuw, ``Silhouettes: A graphical aid to the interpretation and
validation of cluster analysis,'' \emph{J. Comput. Appl. Math.}, vol.~20,
pp.~53--65, 1987, doi: 10.1016/0377-0427(87)90125-7. [Online]. Available:
\url{https://core.ac.uk/download/pdf/82326854.pdf}

\bibitem{arthur2007kmeanspp}
D.~Arthur and S.~Vassilvitskii, ``k-means++: The advantages of careful
seeding,'' in \emph{Proc. 18th Annu. ACM-SIAM Symp. Discrete Algorithms
(SODA)}, 2007, pp.~1027--1035. [Online]. Available (Stanford technical
report): \url{http://ilpubs.stanford.edu:8090/778/}

\bibitem{john2023exploration}
J.~M. John, O.~Shobayo, and B.~Ogunleye, ``An exploration of clustering
algorithms for customer segmentation in the UK retail market,''
\emph{Analytics}, vol.~2, no.~4, pp.~809--823, 2023,
doi: 10.3390/analytics2040042. [Online]. Available:
\url{https://doi.org/10.3390/analytics2040042}

\bibitem{ester1996dbscan}
M.~Ester, H.-P. Kriegel, J.~Sander, and X.~Xu, ``A density-based algorithm
for discovering clusters in large spatial databases with noise,'' in
\emph{Proc. 2nd Int. Conf. Knowledge Discovery and Data Mining (KDD)}, 1996,
pp.~226--231. [Online]. Available:
\url{https://cdn.aaai.org/KDD/1996/KDD96-037.pdf}

\bibitem{ikotun2025cluster}
A.~M. Ikotun, F.~Habyarimana, and A.~E. Ezugwu, ``Cluster validity indices
for automatic clustering: A comprehensive review,'' \emph{Heliyon}, vol.~11,
no.~2, Art.~no.~e41953, 2025, doi: 10.1016/j.heliyon.2025.e41953. [Online].
Available: \url{https://pmc.ncbi.nlm.nih.gov/articles/PMC11787482/}

\bibitem{davidson2005clustering}
I.~Davidson and S.~S. Ravi, ``Clustering with constraints: Feasibility
issues and the k-means algorithm,'' in \emph{Proc. SIAM Int. Conf. Data
Mining (SDM)}, 2005, pp.~138--149. [Online]. Available:
\url{https://web.cs.ucdavis.edu/~davidson/Publications/SDM2005Final.pdf}

\bibitem{gonzalez2025semi}
G.~Gonz\'alez-Almagro, D.~Peralta, E.~De~Poorter, J.-R. Cano, and
S.~Garc\'ia, ``Semi-supervised constrained clustering: An in-depth overview,
ranked taxonomy and future research directions,'' \emph{Artif. Intell.
Rev.}, vol.~58, Art.~no.~157, 2025, doi: 10.1007/s10462-024-11103-8.
[Online]. Available: \url{https://doi.org/10.1007/s10462-024-11103-8}

\bibitem{uci2015onlineretail}
D.~Chen, ``Online Retail,'' UCI Machine Learning Repository, 2015.
[Online]. Available: \url{https://doi.org/10.24432/C5BW33}

\bibitem{pedregosa2011scikit}
F.~Pedregosa \emph{et~al.}, ``Scikit-learn: Machine learning in Python,''
\emph{J. Mach. Learn. Res.}, vol.~12, pp.~2825--2830, 2011. [Online].
Available: \url{https://jmlr.org/papers/v12/pedregosa11a.html}

\bibitem{warrens2022understanding}
M.~J. Warrens and H.~van~der Hoef, ``Understanding the adjusted Rand index and
other partition comparison indices based on counting object pairs,''
\emph{J. Classification}, vol.~39, no.~3, pp.~487--509, 2022,
doi: 10.1007/s00357-022-09413-z. [Online]. Available:
\url{https://link.springer.com/content/pdf/10.1007/s00357-022-09413-z.pdf}

\bibitem{lakens2013calculating}
D.~Lakens, ``Calculating and reporting effect sizes to facilitate cumulative
science: A practical primer for t-tests and ANOVAs,'' \emph{Front. Psychol.},
vol.~4, Art.~no.~863, 2013, doi: 10.3389/fpsyg.2013.00863. [Online].
Available: \url{https://pmc.ncbi.nlm.nih.gov/articles/PMC3840331/}

\bibitem{fielperes2026effect}
F.~Fiel Peres, ``Effect sizes for nonparametric tests,'' \emph{Biochem.
Med. (Zagreb)}, vol.~36, no.~1, Art.~no.~010101, 2026,
doi: 10.11613/BM.2026.010101. [Online]. Available:
\url{https://pmc.ncbi.nlm.nih.gov/articles/PMC12701665/}

\bibitem{virtanen2020scipy}
P.~Virtanen \emph{et~al.}, ``SciPy 1.0: Fundamental algorithms for scientific
computing in Python,'' \emph{Nature Methods}, vol.~17, no.~3, pp.~261--272,
2020, doi: 10.1038/s41592-019-0686-2. [Online]. Available:
\url{https://pmc.ncbi.nlm.nih.gov/articles/PMC7056644/}

\bibitem{rousselet2021percentile}
G.~A. Rousselet, C.~R. Pernet, and R.~R. Wilcox, ``The percentile bootstrap:
A primer with step-by-step instructions in R,'' \emph{Adv. Methods Pract.
Psychol. Sci.}, vol.~4, no.~1, 2021, doi: 10.1177/2515245920911881.
[Online]. Available: \url{https://doi.org/10.1177/2515245920911881}

\bibitem{jolliffe2016pca}
I.~T. Jolliffe and J.~Cadima, ``Principal component analysis: A review and
recent developments,'' \emph{Phil. Trans. R. Soc. A}, vol.~374, no.~2065,
Art.~no.~20150202, 2016, doi: 10.1098/rsta.2015.0202. [Online]. Available:
\url{https://doi.org/10.1098/rsta.2015.0202}

\bibitem{vandermaaten2008tsne}
L.~van~der Maaten and G.~Hinton, ``Visualizing data using t-SNE,''
\emph{J. Mach. Learn. Res.}, vol.~9, pp.~2579--2605, 2008. [Online].
Available: \url{https://jmlr.org/papers/v9/vandermaaten08a.html}

\end{thebibliography}

\end{document}
